\chapter{Artefacts}
\label{app:artefacts}

Every technical chapter of this thesis is accompanied by a mechanized artefact.
All of them are open source, and all are developed in the open rather than
released as a snapshot at publication time. This appendix collects them, so that
a reader who wants to run the procedures rather than read about them knows where
to look.

\paragraph{\Cref{ch:lean-mlir}: the \messa{X} framework.}
The framework is developed openly on GitHub~\cite{leanmlirgithub}. The artefact
accompanying the ITP~2024 publication~\citep{leanmlir} is archived on Zenodo at
\url{https://zenodo.org/records/11519739}.

\paragraph{\Cref{ch:mono-width}: mono-width parametric bitvectors.}
The data supporting this chapter are openly available on Zenodo at
\url{https://doi.org/10.5281/zenodo.16269885}~\cite{bhat_2025_16269885}. This
artefact has been independently reproduced through the published artifact
evaluation of OOPSLA~2025~\citep{bhat2025certified}; it is the only artefact in
this thesis to have been through that process at the time of writing.

\paragraph{\Cref{ch:multi-width}: multi-width parametric bitvectors.}
The evaluation is packaged into a Docker container as the artefact for the
multi-width results~\citep{bhat2026multiwidth}.

\paragraph{\Cref{ch:floating-point}: floating-point bitblasting.}
The mechanization lives at \url{https://github.com/opencompl/fp.lean}. Throughout
\cref{ch:floating-point}, individual definitions and theorems are hyperlinked
directly into this repository, so that each claim in the prose can be checked
against the Lean source that discharges it.
